\documentclass[12pt]{article}
\usepackage[utf8]{inputenc}
\usepackage[T1]{fontenc}
\usepackage[margin=1in]{geometry}
\usepackage{setspace}
\usepackage{xcolor}
\usepackage[colorlinks=true, urlcolor=blue!60!black, linkcolor=blue!60!black]{hyperref}
\setlength{\parindent}{2em}
\setlength{\parskip}{1.25em}
\renewcommand{\baselinestretch}{1.0}
\title{\textit{Cooper v. Aaron}, 358 U.S. 1 (1958)}
\author{Case Brief}
\date{}

\begin{document}
\maketitle

\section{Case Metadata}
\begin{description}
  \item[Citation:] \textit{Cooper v. Aaron}, 358 U.S. 1 (1958) (August Special Term, 1958).
  \item[Argued:] August 28, 1958 (motions) \& September 11, 1958 (merits).
  \item[Decided:] September 12, 1958 (unanimous per curiam judgment); full opinion announced September 29, 1958.
  \item[Vote:] 9--0 (unanimous).
  \item[Author:] Jointly signed by all nine Justices --- Chief Justice Earl Warren, and Justices Hugo L. Black, Felix Frankfurter, William O. Douglas, Harold H. Burton, Tom C. Clark, John M. Harlan II, William J. Brennan, Jr., and Charles E. Whittaker (the only time in Supreme Court history that every Justice personally signed the opinion of the Court).
  \item[Concurrences / Dissents:] Concurring opinion by Justice Felix Frankfurter (358 U.S. at 20--26). No dissents.
  \item[Chapter / Topic:] Judicial Review and the Federal Courts --- Judicial Supremacy, the Supremacy Clause, and Binding Effect of Supreme Court Constitutional Interpretations (U.S. Const.\ art.\ VI, cl.~2--3; amend.\ XIV).
  \item[Primary Source URL:] \href{https://supreme.justia.com/cases/federal/us/358/1/}{Justia U.S.\ Supreme Court Center} (\url{https://supreme.justia.com/cases/federal/us/358/1/}).
\end{description}

\section{Facts and Procedural History}
Three days after \textit{Brown v. Board of Education} (\textit{Brown I}), 347 U.S. 483 (1954), the Little Rock, Arkansas School Board announced its intention to comply with the Supreme Court's mandate. In May 1955---one week before \textit{Brown II}, 349 U.S. 294 (1955)---the Board adopted a gradual desegregation plan scheduled to begin at Central High School in fall 1957 and achieve full district-wide desegregation by 1963. Federal courts upheld the plan against a challenge by Black plaintiffs seeking faster integration (\textit{Aaron v. Cooper}, 143 F.\ Supp.\ 855 (E.D.\ Ark.\ 1956), \textit{aff'd}, 243 F.2d 361 (8th Cir.\ 1957)).

Meanwhile, Arkansas state officials launched a campaign of ``Massive Resistance.'' In November 1956, Arkansas voters adopted Amendment 44 to the State Constitution commanding the General Assembly to oppose ``in every Constitutional manner the Unconstitutional desegregation decisions'' of the U.S.\ Supreme Court, followed in February 1957 by statutes relieving children from compulsory attendance at integrated schools and creating a State Sovereignty Commission.

On September 2, 1957---the night before nine Black students (the ``Little Rock Nine'') were scheduled to enter Central High School---Governor Orval Faubus, without any request from local authorities or any prior crowd violence, deployed Arkansas National Guard units to place the school ``off limits'' to Black students. After the District Court enjoined the Governor and Guard officers on September 20, 1957 (\textit{Aaron v. Cooper}, 156 F.\ Supp.\ 220, \textit{aff'd sub nom.\ Faubus v. United States}, 254 F.2d 797 (8th Cir.\ 1958)), the Guard withdrew, but a hostile mob forced police to remove the students on September 23. President Dwight D. Eisenhower then dispatched regular U.S.\ Army troops (later replaced by federalized National Guardsmen) to enforce the court's orders, allowing eight of the nine students to complete the 1957--1958 academic year.

On February 20, 1958, citing ``chaos, bedlam and turmoil'' and severe disruption to the educational program, the Little Rock School Board and Superintendent petitioned the District Court to suspend their desegregation plan for two and one-half years and return the Black students to segregated schools. On June 20, 1958, the District Court granted the suspension (163 F.\ Supp.\ 13). On August 18, 1958, the Eighth Circuit reversed (257 F.2d 33) but stayed its mandate to permit certiorari. The Supreme Court convened in an extraordinary August Special Term, granted certiorari on September 11, 1958, issued a unanimous per curiam affirmance on September 12 (three days before school was scheduled to open), and released its full jointly signed opinion on September 29, 1958.

\section{Issue(s) Presented}
\begin{enumerate}
  \item May a local school board acting in good faith obtain a judicial suspension of a court-approved desegregation plan because of public hostility, disorder, and violence directly provoked by the official resistance of the State's Governor and Legislature?
  \item Are state executive, legislative, and judicial officials constitutionally bound by the Supreme Court's interpretation of the Fourteenth Amendment in \textit{Brown v. Board of Education} under Article VI of the Constitution?
\end{enumerate}

\section{Holding}
No to the first question; yes to the second. Judgment of the Eighth Circuit \textbf{affirmed}, and District Court orders enforcing the desegregation plan \textbf{reinstated effective immediately}. Local school officials act as agents of the State under the Fourteenth Amendment and cannot invoke their own good faith---or disorder generated by other state officials---to suspend constitutional rights. Moreover, under \textit{Marbury v. Madison} and Article VI's Supremacy and Oath Clauses, the federal judiciary is supreme in the exposition of the law of the Constitution, making the Supreme Court's interpretation of the Fourteenth Amendment in \textit{Brown} the ``supreme Law of the Land'' binding on all state officials.

\section{Reasoning}
The Court's unanimous, jointly signed opinion rests on two core pillars:
\begin{enumerate}
  \item \textbf{State Action and Rejection of the ``Heckler's Veto'' of Official Disorder:} The Court accepted both the School Board's good faith and the District Court's findings that the 1957--1958 school year had been marked by turmoil. Nevertheless, the record indisputably showed that those conditions were ``directly traceable'' to the official actions of the Arkansas Governor and Legislature. Under \textit{Ex parte Virginia}, 100 U.S. 339, 347 (1879), a State acts through all of its agencies and officers; for Fourteenth Amendment purposes, the School Board stood as an agent of the State and could not plead obstruction by other state officials as a legal excuse. Invoking \textit{Buchanan v. Warley}, 245 U.S. 60, 81 (1917), the Court emphasized that ``law and order are not here to be preserved by depriving the Negro children of their constitutional rights.'' Rights declared in \textit{Brown} cannot be nullified openly or through evasive schemes, ``whether attempted `ingeniously or ingenuously''' (\textit{Smith v. Texas}, 311 U.S. 128, 132 (1940)).
  \item \textbf{Judicial Supremacy under Article VI and \textit{Marbury}:} Although the state-action holding sufficed to dispose of the case, the Court deliberately went further to answer the Arkansas Governor and Legislature's premise that they were not bound by \textit{Brown} because they were not parties to it. Reading Article VI's Supremacy Clause together with Chief Justice Marshall's declaration in \href{https://supreme.justia.com/cases/federal/us/5/137/}{\textit{Marbury v. Madison}}, 5 U.S. (1 Cranch) 137, 177 (1803) (``It is emphatically the province and duty of the judicial department to say what the law is''), the Court reasoned that \textit{Marbury} ``declared the basic principle that the federal judiciary is supreme in the exposition of the law of the Constitution.'' Consequently, the Court's \emph{interpretation} of the Fourteenth Amendment in \textit{Brown} is itself the ``supreme law of the land,'' binding state legislators, executives, and judges under their Article VI, Clause 3 oath ``to support this Constitution.'' Citing \textit{Ableman v. Booth}, 62 U.S. (21 How.) 506 (1859), \textit{United States v. Peters}, 9 U.S. (5 Cranch) 115 (1809), and \textit{Sterling v. Constantin}, 287 U.S. 378 (1932), the Court concluded that no state official ``can war against the Constitution without violating his undertaking to support it,'' and noted that the three Justices appointed since \textit{Brown} unanimously reaffirmed it.
\end{enumerate}

\section{Concurring and Dissenting Opinions}
\begin{itemize}
  \item \textbf{Justice Frankfurter (Concurring, 358 U.S. at 20--26):} Writing separately while ``unreservedly participating'' in the joint opinion (aimed largely at Southern lawyers and leaders), Frankfurter stressed that peaceful community accommodation had been underway in Little Rock before Arkansas deployed the militia to thwart law. Yielding to state-sponsored violence would ``enthrone official lawlessness'' and signify ``that law should bow to force.'' Quoting his own concurrence in \textit{United States v. United Mine Workers}, 330 U.S. 258, 307--309 (1947), Frankfurter distinguished permissible criticism and orderly legal revision (reconsideration, legislation, or constitutional amendment) from active defiance, closing with an appeal to Lincoln's ``better angels of our nature.''
\end{itemize}

\section{Key Quotations}
\begin{quote}
``The constitutional rights of respondents are not to be sacrificed or yielded to the violence and disorder which have followed upon the actions of the Governor and Legislature.'' (358 U.S. at 16)
\end{quote}

\begin{quote}
``In short, the constitutional rights of children not to be discriminated against in school admission on grounds of race or color declared by this Court in the Brown case can neither be nullified openly and directly by state legislators or state executive or judicial officers nor nullified indirectly by them through evasive schemes for segregation whether attempted `ingeniously or ingenuously.''' (358 U.S. at 17)
\end{quote}

\begin{quote}
``This decision declared the basic principle that the federal judiciary is supreme in the exposition of the law of the Constitution, and that principle has ever since been respected by this Court and the Country as a permanent and indispensable feature of our constitutional system. It follows that the interpretation of the Fourteenth Amendment enunciated by this Court in the Brown case is the supreme law of the land, and Art.\ VI of the Constitution makes it of binding effect on the States `any Thing in the Constitution or Laws of any State to the Contrary notwithstanding.''' (358 U.S. at 18)
\end{quote}

\begin{quote}
``No state legislator or executive or judicial officer can war against the Constitution without violating his undertaking to support it.'' (358 U.S. at 18)
\end{quote}

\begin{quote}
``Violent resistance to law cannot be made a legal reason for its suspension without loosening the fabric of our society. What could this mean but to acknowledge that disorder under the aegis of a State has moral superiority over the law of the Constitution?'' (358 U.S. at 22, Frankfurter, J., concurring)
\end{quote}

\section{Doctrinal Significance and Modern Legacy (Through 2026)}
\textit{Cooper v. Aaron} is the canonical judicial statement of \textbf{judicial supremacy}---the claim not merely that federal courts have the power of judicial review in litigated cases (\textit{Marbury}), nor merely that specific court judgments bind parties, but that the Supreme Court's constitutional \emph{interpretations} themselves constitute the ``supreme Law of the Land'' binding non-party officials nationwide.
\begin{itemize}
  \item \textbf{Expansion to Coordinate Federal Branches:} Although \textit{Cooper} involved state defiance (vertical federalism), the Supreme Court subsequently extended \textit{Cooper}'s judicial-supremacy framework horizontally against Congress and the President---most notably in \href{https://supreme.justia.com/cases/federal/us/418/683/}{\textit{United States v. Nixon}}, 418 U.S. 683, 703--705 (1974) (invoking \textit{Marbury} to reject presidential finality over executive privilege), and \href{https://supreme.justia.com/cases/federal/us/521/507/}{\textit{City of Boerne v. Flores}}, 521 U.S. 507, 535--536 (1997) (striking down the Religious Freedom Restoration Act as applied to states because Congress may not alter the Court's substantive interpretation of the Fourteenth Amendment).
  \item \textbf{Judicial Supremacy vs.\ Departmentalism Through 2026:} \textit{Cooper}'s maximalist reading of Article VI has remained contested in both scholarship (see \href{https://pennlawreview.com/2024/03/01/clarifying-judicial-aggrandizement/}{\textit{University of Pennsylvania Law Review} (2024)}) and recent litigation. In \href{https://supreme.justia.com/cases/federal/us/606/24a884/}{\textit{Trump v. CASA, Inc.}}, 606 U.S.\ \_\_\_ (2025) (\href{https://www.supremecourt.gov/opinions/24pdf/24a884_8n59.pdf}{slip op.}), the Supreme Court curtailed district courts' authority to issue ``universal injunctions'' extending relief to non-parties, reigniting debates over the boundary between a binding judgment against parties and the broader precedential force of constitutional rulings---and prompting intense 2025--2026 scrutiny of how the Executive Branch and lower courts must treat Supreme Court interim orders and merits rulings (see \href{https://harvardlawreview.org/wp-content/uploads/2025/11/139-Harv.-L.-Rev.-86.pdf}{139 Harv.\ L.\ Rev.\ 86 (2025)}; \href{https://supreme.justia.com/cases/federal/us/609/25-365/}{\textit{Trump v. Barbara}}, 609 U.S.\ \_\_\_ (2026)).
\end{itemize}

\section{Pedagogical Critiques and Exam Hypotheticals}
\subsection{Key Critiques}
\begin{itemize}
  \item \textbf{Textual Slippage from \textit{Marbury} and Article VI:} Article VI lists three items as ``the supreme Law of the Land'': the Constitution, federal laws made in pursuance thereof, and treaties. It does \emph{not} list Supreme Court opinions. Critics (including former Attorney General Edwin Meese III and scholars of ``departmentalism'' or ``coordinate review'') argue that \textit{Cooper} conflates the Constitution itself with the Court's fallible \emph{interpretations} of it, and leaps from \textit{Marbury}'s modest case-specific judicial review to universal interpretive supremacy.
  \item \textbf{The Lincoln / \textit{Dred Scott} Objection:} If \textit{Cooper} were taken literally, President Abraham Lincoln would have violated his oath when he refused to treat \textit{Dred Scott v. Sandford} as a permanent political rule for Congress and the Executive beyond the specific parties before the Court. Even Justice Frankfurter's concurrence quietly tempered the joint opinion by conceding that ``[e]ven this Court has the last say only for a time'' and may be asked to reconsider its decisions (358 U.S. at 23).
\end{itemize}

\subsection{Exam Hypothetical}
After the Supreme Court holds in \textit{State A v. Doe} that a state ban on a particular medical procedure violates the Fourteenth Amendment, the Governor of State B (not a party to \textit{Doe}) signs an executive order directing state prosecutors to continue enforcing State B's identical statute, declaring: ``My oath is to the Constitution, not to five Justices' mistaken reading of it; I will obey a federal court injunction directed at me, but until then I am not bound by \textit{Doe}.''
\begin{itemize}
  \item \textbf{Analysis:} Under \textit{Cooper v. Aaron}'s strong judicial-supremacy framework, the Governor's refusal to treat the Court's Fourteenth Amendment holding as binding law violates Article VI and the Oath Clause; moreover, in any federal suit against State B's officials, lower federal courts are bound by vertical \textit{stare decisis} to enjoin enforcement immediately (and state officials face potential fee-shifting under 42 U.S.C.\ \S~1988 or loss of qualified immunity because the right is clearly established). Under a \textbf{departmentalist} view (invoking Lincoln's response to \textit{Dred Scott}), the Governor distinguishes between defying a direct court judgment (unlawful) and declining to treat a precedent involving other parties as automatically settling his independent constitutional judgment while seeking a test case to ask the Court to reconsider \textit{Doe}.
  \item \textbf{Opposite Side Contrast (Departmentalism / Lincoln's First Inaugural \& \textit{Marbury} Narrowly Read):} Contrast \textit{Cooper}'s judicial supremacy with classical \textbf{departmentalism} (associated with Jefferson, Jackson, and Lincoln). Whereas \textit{Cooper} treats the Supreme Court's constitutional exposition as identical to the ``supreme Law of the Land'' binding every official nationwide, departmentalists distinguish \emph{judgments} (which bind the parties to a case and must be obeyed) from \emph{constitutional reasoning} (where coordinate branches---and arguably states seeking orderly reconsideration---retain independent interpretive duties, subject to judicial check when a concrete Article III case arises). On an exam, always separate (1) defying an actual court order directed at an official (unlawful under \emph{both} \textit{Cooper} and departmentalism) from (2) taking an independent constitutional position outside a binding decree (barred by \textit{Cooper}'s broad dicta, defended by departmentalists).
\end{itemize}

\end{document}
